\def\SECTION{2.2 Conditional Statements}
\documentclass[aspectratio=169,14pt]{beamer}
\usepackage[T1]{fontenc}
\usepackage{mathpazo}
\usepackage{amsmath,amssymb}
\usepackage{tikz}
\usetikzlibrary{arrows.meta,calc,decorations.pathreplacing}
\definecolor{fblue}{HTML}{1F4E79}
\definecolor{fteal}{HTML}{0F7C7C}
\definecolor{fgold}{HTML}{C9962B}
\definecolor{fred}{HTML}{B3261E}
\definecolor{fink}{HTML}{16233B}
\definecolor{fsoft}{HTML}{EEF3F8}
\definecolor{fgoldsoft}{HTML}{FBF1DC}
\setbeamertemplate{navigation symbols}{}
\setbeamercolor{normal text}{fg=fink}
\setbeamercolor{frametitle}{fg=fblue}
\setbeamerfont{frametitle}{series=\bfseries,size=\Large}
\setbeamercolor{structure}{fg=fteal}
\setbeamertemplate{itemize item}{\color{fteal}$\blacktriangleright$}
\setbeamertemplate{frametitle}{\vspace{6pt}\insertframetitle\par\vspace{-2pt}{\color{fgold}\rule{2.2cm}{2pt}}}
\setbeamertemplate{footline}{\hspace{0.6cm}{\color{fblue}\scriptsize\bfseries Fang\textcolor{fteal}{Math}}\hfill{\scriptsize\color{gray}\SECTION\quad\insertframenumber}\hspace{0.6cm}\vspace{6pt}}
\newcommand{\key}[1]{{\color{fteal}\textbf{#1}}}
\newcommand{\ans}[1]{\colorbox{fgoldsoft}{\color{fink}\ensuremath{#1}}}
\newcommand{\warn}[1]{\begin{beamercolorbox}[sep=8pt,rounded=true]{warnbox}{\color{fred}\textbf{Watch out!}}\enskip #1\end{beamercolorbox}}
\setbeamercolor{warnbox}{bg=fred!8,fg=fink}
\setbeamercolor{turnbox}{bg=fgoldsoft,fg=fink}
\setbeamercolor{defbox}{bg=fsoft,fg=fink}
\newcommand{\defn}[1]{\begin{beamercolorbox}[sep=8pt,rounded=true]{defbox}#1\end{beamercolorbox}}
\newenvironment{yourturn}{\begin{beamercolorbox}[sep=10pt,rounded=true]{turnbox}{\color{fgold}\Large\bfseries Your turn!}\quad{\small Pause the video and try it.}\par\medskip}{\end{beamercolorbox}}
\newcommand{\titleframe}[3]{%
\begin{frame}[plain]
\begin{tikzpicture}[remember picture,overlay]
\fill[fblue] (current page.north west) rectangle ([yshift=-0.9cm]current page.north east);
\node[anchor=west,text=white,font=\bfseries] at ([xshift=0.6cm,yshift=-0.45cm]current page.north west) {Fang\textcolor{fgold}{Math}};
\node[anchor=east,text=white,font=\small] at ([xshift=-0.6cm,yshift=-0.45cm]current page.north east) {Geometry $\cdot$ Chapter \CHNO};
\end{tikzpicture}
\vfill
{\color{fteal}\large\bfseries #1}\par\medskip
{\color{fblue}\fontsize{30}{34}\selectfont\bfseries #2}\par\medskip
{\color{fgold}\rule{3cm}{3pt}}\par\bigskip
{\large #3}
\vfill
{\small\color{gray} fangmath.com $\cdot$ YouTube @fangmathcom}
\end{frame}}
\newcommand{\closeframe}[1]{%
\begin{frame}{Remember}
#1
\vfill
\begin{center}{\color{fblue}\large\bfseries Now practice at your level:}\par\smallskip
{\Large\color{fteal}\bfseries fangmath.com} \enskip{\color{gray}$\to$ Geometry $\to$ Chapter \CHNO}\end{center}
\end{frame}}

\usepackage{pgfplots}\pgfplotsset{compat=1.17}
\newcommand{\grid}[4]{\draw[step=1,gray!25] (#1,#2) grid (#3,#4);\draw[-{Stealth},gray] (#1,0)--(#3+0.3,0) node[right]{\small$x$};\draw[-{Stealth},gray] (0,#2)--(0,#4+0.3) node[above]{\small$y$};}
\def\CHNO{2}
\usetikzlibrary{angles,quotes}
\begin{document}
\titleframe{Lesson 2.2}{Conditional Statements}{If\ldots\ then\ldots --- and its converse, inverse and contrapositive}
\begin{frame}{If $p$, then $q$}
\onslide<1->{{\large If \colorbox{fteal!15}{it rains}, then \colorbox{fgold!25}{the ground is wet}.}}\par\medskip
\onslide<2->{{\color{fteal}\textbf{$p$ = hypothesis}} (the ``if'' part) \qquad {\color{fgold}\textbf{$q$ = conclusion}} (the ``then'' part) \qquad written $p \to q$}\par\medskip
\end{frame}
\begin{frame}{Four related statements}
\onslide<1->{\key{Conditional} \quad $p \to q$ \quad If it rains, then the ground is wet.}\par\medskip
\onslide<2->{\key{Converse} \quad $q \to p$ \quad \textbf{switch}: If the ground is wet, then it rained.}\par\medskip
\onslide<3->{\key{Inverse} \quad $\sim p \to \sim q$ \quad \textbf{negate}: If it does not rain, then the ground is not wet.}\par\medskip
\onslide<4->{\key{Contrapositive} \quad $\sim q \to \sim p$ \quad \textbf{switch and negate}: If the ground is not wet, then it did not rain.}\par\medskip
\end{frame}
\begin{frame}{True or false?}
\onslide<1->{Conditional: If it rains, the ground is wet. \quad {\color{fteal}\textbf{true}}}\par\medskip
\onslide<2->{Converse: If the ground is wet, it rained. \quad {\color{fred}\textbf{false}} --- a sprinkler!}\par\medskip
\onslide<3->{Contrapositive: If the ground is not wet, it did not rain. \quad {\color{fteal}\textbf{true}}}\par\medskip
\end{frame}
\begin{frame}{The rule}
\onslide<1->{\defn{A conditional and its \key{contrapositive} are always both true or both false. \\ The \key{converse} and the \key{inverse} also match each other --- but they may not match the original.}}\par\medskip
\end{frame}
\begin{frame}{Biconditional}
\onslide<1->{If a conditional \textbf{and} its converse are both true, write ``\key{if and only if}''.}\par\medskip
\onslide<2->{An angle is right \key{if and only if} it measures $90^\circ$.}\par\medskip
\end{frame}
\begin{frame}{}
\begin{yourturn}``If an angle measures $30^\circ$, then it is acute.'' \\ Write the \textbf{converse}. Is it true?\end{yourturn}\par\medskip
\onslide<2->{Converse: If an angle is acute, then it measures $30^\circ$. \quad \ans{\text{false}} \quad counterexample: a $50^\circ$ angle is acute.}\par\medskip
\end{frame}
\closeframe{\begin{itemize}\setlength\itemsep{8pt}
\item Conditional $p \to q$: \key{converse} switches, \key{inverse} negates, \key{contrapositive} does both.
\item A conditional and its contrapositive are both true or both false.
\item Both directions true? Use ``\key{if and only if}''.
\end{itemize}}
\end{document}
